\documentclass[11pt,a4paper]{article}
\usepackage[T1]{fontenc}
\usepackage[utf8]{inputenc}
\usepackage{amsmath,amssymb,xcolor}
\usepackage[margin=25mm]{geometry}
\usepackage[colorlinks=true,urlcolor=blue]{hyperref}
\definecolor{nexorablue}{HTML}{2459B5}
\setlength{\parindent}{0pt}
\setlength{\parskip}{7pt}
\setlength{\emergencystretch}{2em}
\renewcommand{\familydefault}{\sfdefault}
\begin{document}

{\small\bfseries\color{nexorablue} NEXORA CPGE}\par
{\small MP / NEX-SLC-06 / Corrigé}\par
{\LARGE\bfseries Choisir le plus petit gain conforme au cahier des charges}\par
\vspace{0.25em}\hrule\vspace{0.6em}

\section*{Deux inégalités}
Le transfert est $K/(p+4+K)$ et son pôle est négatif. La constante de temps $1/(4+K)$ impose $K\ge6$. L’écart permanent $1-K/(4+K)=4/(4+K)$ impose $K\ge16$.

\section*{Choix final}
Le plus petit gain compatible est $16$. Il donne $\tau_f=0{,}050\,\mathrm s$ et un écart de $0{,}20$. La précision est la contrainte active.


\end{document}
